\PassOptionsToPackage{dvipsnames}{xcolor}
%\documentclass[12pt]{beamer}
%\documentclass[10pt]{beamer}
\documentclass[10pt, handout]{beamer}

\usepackage{mathtools} %% previously \usepackage{amsmath}
\usepackage{amsfonts}
\usepackage{amssymb} % http://www.ctan.org/tex-archive/macros/latex/required/psnfss/psnfss2e.pdf
%\usepackage[official]{eurosym}
\usepackage{helvet}
\usepackage{natbib}
\usepackage{lscape, longtable, array, booktabs, multirow, tabularx, graphics, epstopdf}
\usepackage{setspace}
\usepackage{listings}
\usepackage {tikz}

\usefonttheme[onlymath]{serif}

\usepackage{accents}
\newcommand*{\dt}[1]{%
   \accentset{\mbox{\LARGE $.$}}{#1}}

%\usepackage[T1]{fontenc}

\usecolortheme{dolphin}
%\useoutertheme{infolines}
\definecolor{dollarbill}{rgb}{0.52, 0.73, 0.4}
%\setbeamercovered{transparent}% Dim out "inactive" elements
\setbeamercovered{invisible}% Default
\beamertemplatenavigationsymbolsempty
\setbeamerfont{page number in head/foot}{size=\footnotesize}
\setbeamertemplate{footline}[frame number]


\usepackage{hyperref} 
\hypersetup{colorlinks=true}
%\hypersetup{linkcolor=RoyalBlue, urlcolor=ForestGreen, citecolor=Magenta}
\hypersetup{allcolors=RoyalBlue} %% 

%\renewcommand{\thefootnote}{\fnsymbol{footnote}}
\renewcommand{\thefootnote}{\arabic{footnote}}

\DeclareMathOperator*{\argmax}{argmax} % thin space, limits underneath in displays


\title [Recursive Methods]{Recursive Methods}
\vskip 0.1in

\subtitle{Mathematical Methods for Economics (771)}
%\author[H Hollander]{Hylton Hollander}

\institute[Stellenbosch University]{\large{
Stellenbosch University}}


%\date[April, 2020]{April, 2020}

\date{}

\begin{document}

\frame{\titlepage}

\frame{\frametitle{Contents:} \tableofcontents}


\begin{frame}

\section{Readings}

\frametitle{Readings} \vskip 0.1in

\begin{itemize}
\item Krusell, Per. (2014). Real Macroeconomic Theory, Chapter 4 ``Dynamic Optimization''

\item Judd, Kenneth, L. (1991). ``A Review of \textit{Recursive Methods in Economic Dynamics} by Stockey, N., Lucas Jr., R., and Prescott, E.,'' Journal of Economic Literature, 29(1), 69$-$77
\end{itemize}

Additional:
\begin{itemize}
\item  Namay L. Stockey, Robert E. Lucas Jr., and Edward C. Prescott, 1989. Recursive methods in economic dynamics. Harvard university
press (CH 4, 5).
\end{itemize}


\end{frame}


\section{Introduction}
\subsection{Why dynamic programming?}


\begin{frame}

\frametitle{Why Dynamic Programming?} \vskip 0.1in
% ref: McCandless: 50

%$\square$ Recursive methods (dynamic programming) is a fundamental tool of dynamic economic analysis.

$\square$ Most macro models involve a dynamic optimization problem and a resulting (infinite) sequence of real numbers that solves it.
\begin{itemize}
\item Dynamic decisions are made recursively (time period by time period) and not once-and-for-all. For example, savings between $t$ and $t+1$ are decided on at $t$, and not $0$

\item But: the nature/structure of the optimization problem that a decision maker faces does not depend on the period in which they are making their decisions: it is identical (stationary) at every point in time. 
%\item the nature of the problems that we are going to solve do not change with time;
\item What changes from period to period are the initial conditions
$\rightarrow$ the values of the variables that have been determined
by the past or by ``nature''.
\end{itemize}

$\square$ The search for a sequence is sometimes impractical, and not always intuitive. An alternative approach is often available: \textit{dynamic programming}  


\end{frame}


\begin{frame}

\frametitle{Why Dynamic Programming?} \vskip 0.1in
%see Krussel p.58

$\square$ Recursive methods (dynamic programming) is a fundamental tool of dynamic economic analysis:
\begin{itemize}
\item Useful conceptually as well as for analytical and, especially, \textit{numerical} computation.

\item Allows for comparative dynamic exercises, in dynamic models with uncertainty (e.g., a counterfactual policy intervention)

\item Permit the inclusion of the stochastic shocks (i.e., a non-deterministic system).

\item But the reduction of a dynamic model to a \textit{recursive model} must be done carefully. 

\end{itemize}

$\square$ We will go over the basics of this approach.

$\square$ The focus will be on concepts, as opposed to mathematical rigor or formal proofs


%see Krussel p.58
% The nature of the optimization problem that individuals face does not depend on the period in which they are making their decisions
\end{frame}



%\subsection{Recursive methods in economic dynamics}

\begin{frame}

\frametitle{Recursive methods in economic dynamics} %\vskip 0.1in
% ref: McCandless: 50
% for further discussion of the inclusion of stochastic shocks see McCandless CH5

\begin{quote}
In general, we turn away from looking for a sequence of prices and allocations that satisfy equilibrium conditions, and instead look for a collection of policy functions, independent of time, which express current decisions and prices as functions of the state variables, which in turn are sufficient statistics of the past. 
\begin{flushright}
--- Judd (1991, p.71)
\end{flushright} 
\end{quote} \vskip 0.2in

\end{frame}



%\subsection{States and controls}

\begin{frame}

\frametitle{Preliminaries} \vskip 0.1in


%The reduction of a dynamic model to a recursive model must be done carefully. 
$\square$  We are looking for a function $g(\cdot)$ that does not vary with time: a \textit{decision} or \textit{policy rule} \vskip0.1in


$\square$ The critical step is defining the \textcolor{blue}{state variable} 
%\textcolor{blue}{States and Controls}
\begin{itemize}
\item variables whose values are already determined in period $t$ (predetermined)
\end{itemize} \vskip0.1in
%\pause

$\square$ The choice of \textcolor{blue}{control variables} can matter in how easily we can solve the model. 
\begin{itemize}
\item decision variables whose values decision makers explicitly choose in period $t$ with the goal of optimizing their objective function.
\end{itemize} \vskip0.1in

$\square$ Whatever the choice, there must be enough constraints (market conditions) such that the values of the rest of the relevant variables in period $t$ are determined.

\end{frame}




\section{A simple problem and an alternative approach to its solution}


\subsection{The canonical problem}

\begin{frame}
\frametitle{ The canonical problem}
We consider the finite horizon dynamic optimization problem:%\footnote{}
\begin{eqnarray}\label{Eq1}
\max_{\{k_{t+1}\}^{T}_{t=0}} && \sum^{T}_{t=0} \beta^t F(k_t, k_{t+1}) \\ \nonumber
\text{s.t.} && k_{t+1} \geq 0 ~ \forall ~ t
\end{eqnarray}
where we want to find the sequence $\{k_{t+1}\}^{T}_{t+0}$ that maximizes the objective function (\ref{Eq1}). \vskip0.25in

Assume $F(k_t, k_{t+1}) = u(c_t)$ is a concave utility function, and $\beta$ is the stationary discount factor. Let $k_t$ be the capital stock at \textit{the beginning of period} $t$, $y_t = f(k_t)$ a neoclassical production function, and $c_t$ consumption in period $t$ chosen at \textit{the end of the period}.

%\begin{equation}
%u(c_t)
%\end{equation}
% additively separable, with stationary discounting weights.



\end{frame}




\begin{frame}
\frametitle{ The canonical problem}
Then a social planner for this economy will solve the problem
\begin{eqnarray}\label{Eq2}
\max_{\{c_{t}\}} \sum^{T}_{t=0} \beta^t u(c_{t})
\end{eqnarray}
where the choice path for $c_t$ is constrained by production possibilities, represented by the law of motion
\begin{equation}\label{Eq3}
k_{t+1} = \underset{\text{Investment}}{\underbrace{y_t - c_t}} = f(k_t) - c_t,
\end{equation} 
where $k_0 > 0$ is the initial endowment of capital. \vskip0.15in

To make our lives simple we assume:
\begin{itemize} \scriptsize
\item the resource constraint (\ref{Eq3}) is binding; 
\item $c_t$ and $k_t$ are nonnegative for all $t$; 
\item $\underset{{c \to 0}}{\lim} u'(c)=\infty$ (i.e., $c_t = 0$ at any $t$ cannot be optimal). So we ignore $c_{t} \geq 0$, but we still consider the inequality $k_{t+1} \geq 0$ (Krussel, p.44); and, 
\item $\delta = 1$ implies that the capital stock depreciates 100\%. 
\end{itemize}

\end{frame}


\begin{frame}
\frametitle{ The canonical problem}
We now have a consumption-savings decision problem, with the following Lagrangian function:
\begin{equation}
\mathcal{L} = \sum^{T}_{t=0} \beta^t [u[f(k_t) - k_{t+1}] + {\mu}k_{t+1}]
\end{equation}
The next step involves taking the derivative w.r.t the decision variable $k_{t+1}$. The first-order conditions are:
%(recall our assumption about $c_t$?)
\begin{eqnarray}\nonumber
\frac{\partial\mathcal{L}}{\partial{k_{t+1}}} &:& - \beta^{t}u'(c_t) + \beta^{t}\mu_t + \beta^{t+1}u'(c_{t+1})f'(k_{t+1}) = 0, ~~~ t=0, \ldots , T-1 \\ \nonumber
\frac{\partial\mathcal{L}}{\partial{k_{T+1}}} &:& - \beta^{T}u'(c_T) + \beta^{T}\mu_T = 0  , ~~~ {\text{for period}}~T
\end{eqnarray}
Finally, the Kuhn-Tucker conditions include:
\begin{eqnarray}\nonumber
\mu_{t}k_{t+1} = 0, ~~~ t=0, \ldots , T \\ \nonumber
k_{t+1} \geq 0, ~~~ t=0, \ldots , T \\ \nonumber
\mu_{t} \geq 0, ~~~ t=0, \ldots , T 
\end{eqnarray}

\end{frame}

\begin{frame}
\frametitle{ The canonical problem}
The summary statement of the first-order conditions is then the ``Euler equation'':
\begin{eqnarray} \nonumber
u'[f(k_t)-k_{t+1}] &=& {\beta}u'[f(k_{t+1})-k_{t+2}]f'(k_{t+1}) ~, ~~~ t = 0, \ldots, T-1 \\ \nonumber
 &&  k_{0} ~{\text{given}}, ~ k_{T+1}=0 ,
\end{eqnarray}
where the capital sequence is what we need to solve for.\footnote{recall the optimization problem (\ref{Eq1})} \vskip0.1in

The following conditions ensure a unique solution, such that the FOCs are sufficient:\footnote{\scriptsize That is, computing the equilibrium policy function in a recursive model is valuable because it is a sufficient description of equilibrium, and from it one can derive any economic quantity (Judd, 1991)}
\begin{itemize}
\item The objective $U = \sum^{T}_{t=0} u(c_t)$ is (strictly) concave 
\item The constraint (choice) set is convex in $\{c_t , k_{t+1}\}$
\end{itemize}

\vfill

{\scriptsize {$1$}Recall the optimization problem (\ref{Eq1}).}

{\scriptsize $2$ That is, computing the equilibrium policy function in a recursive model is valuable because it is a sufficient description of equilibrium, and from it one can derive any economic quantity (Judd, 1991).}

{\scriptsize Note: From $\partial\mathcal{L}/\partial{k_{T+1}}$, and since $u'(c) > 0 ~ \forall ~ c$, we conclude that $\mu_{T} > 0$. This implies that $k_{T+1} = 0$: the consumer leaves no capital for after the last period.}

\end{frame}

\begin{frame}
\frametitle{ The canonical problem}

Let's interpret the key equation for optimization, the Euler equation:

\begin{equation}
\underset{\text{invest ``one'' more unit}}{\underset{\text{Utility lost if you}}{\underbrace{u'(c_t)}}} 
= 
\underset{\text{next period per unit}}{\underset{\text{Utility increase}}{\underbrace{\beta{u}'(c_{t+1})}}}  \cdot  \underset{\text{invested unit}}{\underset{\text{Return on the }}{\underbrace{f'(k_{t+1})}}}
\end{equation}
Thus, because of the concavity of $u$, equalizing the marginal cost of saving (LHS) to the marginal benefit of saving (RHS) is a condition for an optimum. \vskip0.25in


How do the primitives affect savings behaviour? 

Three components:
\begin{itemize}
\item[(i)] Consumption ``smoothing'' via strictly concave utility function $u$
\item[(ii)] Impatience via discount factor $\beta$
\item[(iii)] Income and substitution effects via the return to savings $f'(k_{t+1}) = R_t$
\end{itemize}

\end{frame}


\begin{frame}
Example 4.1: Logarithmic utility (p. 47)
\end{frame}


\begin{frame}
Example 4.2: CIES (constant intertemporal elasticity of substitution) utility function (p. 49)
\end{frame}




\begin{frame}
\frametitle{Infinite Horizon \& Sufficient conditions}
Why should macroeconomists study the case of an infinite horizon?
\begin{itemize}
\item Altruism: people do not live forever, but they may care about their descendants
\item Simplicity: with a long time horizon, finite- and infinite-horizon models show very similar results. Infinite horizon models are stationary in nature.
\end{itemize}\vskip0.2in


The infinite horizon only requires one additional condition to that in the finite case: the \textit{transversality condition}. Both ensure the capital stock is zero in the limit. See Proposition 4.4, p. 56:



\end{frame}


\subsection{Recursive methods}

\begin{frame}
% A simple problem and an alternative approach to its solution
\frametitle{An alternative approach}
% Krussel, p. 58
%\textcolor{blue}{Dynamic Programming}
\begin{itemize}
\item Our approach up to now has been to look for a sequence of real numbers $\{ k_{t+1} \}^{\infty}_{0}$ that generates an optimal consumption plan.
\item The solution was a difference (functional) equation: the Euler equation.
\item The search for a sequence is sometimes impractical, and not always intuitive.
\item An alternative approach that is intuitive and useful for both analytic \textit{and} numerical computation, is dynamic programming using \textit{recursive methods}
\end{itemize}


\end{frame}





\begin{frame}
% A simple problem and an alternative approach to its solution
\frametitle{The value function}
%\textcolor{blue}{The value function}

Using the canonical (neoclassical) model as an example, assume we can derive the individual's discounted value of utility in period $t$ as:

{ {
\begin{equation}\label{EqA.1}
V(k_t) \equiv \max_{\{k_{t+1+i}\}_{i=0}^{\infty}} \sum_{i=0}^{\infty} \beta^{i} u[f(k_{t+i})- k_{t+1+i}] %~, ~~ {\text s.t.}~ k_{
\end{equation}
%where $\Gamma(k_t)$ represents the feasibale choice set for $k_{t+1}$ given initial condition $k_t$. 
Given the current state ($k_t$), $V(k_t)$ gives the supremum over all possible policies of the present values of current and future utility. 
\vskip0.1in

\textcolor{blue}{The value function} in period ${\color{red}t+1}$:
\begin{equation}\label{EqA.2}
V(k_{t+1}) = \max_{\color{red}\{k_{t+1+i}\}_{i=1}^{\infty}} \sum_{{\color{red}i=1}}^{\infty} \beta^{{\color{red}i-1}} u
[f(k_{t+i})- k_{t+1+i} ]
\end{equation}

We now separate the period $t$ problem \eqref{EqA.1} from that of future periods \ldots

}}
\end{frame}



\begin{frame}
% A simple problem and an alternative approach to its solution

\frametitle{The value function}

\ldots using maximisation-by-steps:
 \begin{eqnarray} \nonumber
V(k_t) & = & \max_{ k_{t+1} \in [0,f(k_t)] } \Big\{ \underset{ i = 0}{\underbrace{u[f(k_{t})- k_{t+1}]}}  \\ \nonumber & & +  \max_{\{k_{t+1+i}\}_{i=1}^{\infty}} \sum_{i=1}^{\infty} \beta^{i} u
[f(k_{t+i})- k_{t+1+i}] \Big \}
 \end{eqnarray}


%\textcolor{blue}{The value function}
\begin{eqnarray}\nonumber
V(k_t) & = & \max_{ k_{t+1} \in [0,f(k_t)] } \Big\{ u[f(k_{t})- k_{t+1}] \\ \label{EqA.3}  & & + ~ \beta \max_{\color{red}\{k_{t+1+i}\}_{i=1}^{\infty}} \sum_{{\color{red}i=1}}^{\infty} \beta^{{\color{red}i-1}} u
[f(k_{t+i})- k_{t+1+i} ] \Big \}
 \end{eqnarray}

By definition of \eqref{EqA.2}, \eqref{EqA.3} equals:

\begin{equation}\nonumber 
V(k_{t}) = \max_{k_{t+1} \in [0,f(k_t)] } \{ u [f(k_{t})-k_{t+1}] +
\beta V(k_{t+1}) \}
\end{equation}


\end{frame}


\begin{frame}

\frametitle{The value function} \vskip 0.1in

%\textcolor{blue}{Remarks:}\\\vskip 0.1in

\begin{equation}\label{EqA.4}
V(k_{t}) = \max_{k_{t+1} \in [0,f(k_t)] } \{ u [f(k_{t})-k_{t+1}] +
\beta V(k_{t+1}) \}
\end{equation}

\begin{itemize}
\item \eqref{EqA.4} is the dynamic programming formulation.
\item It presents exactly the same problem as that shown in \eqref{EqA.1}, but written in a recursive form;
\item It is known as the Bellman equation, and it is a functional equation: the unknown is a function $V$
\item If we find a $V$ that satisfies \eqref{EqA.4} for any value of $k_t$, then all the maximizations on the RHS are well-defined
%\item we denote the value function of the discounted utility by $V(K_{t-1},Z_{t}$), to stress that it is a function of the value of the initial capital stock, $K_{t-1}$ (state variable);
\item The decision rule for $k_{t+1} = g(k_t)$, alluded to earlier, follows:
\end{itemize}
\begin{equation}
g(k_t) = \argmax_{k_{t+1}} \{ u [f(k_{t})-k_{t+1}] +
\beta V(k_{t+1}) \}
\end{equation}
\begin{itemize}
\item To proceed, we need to assume that the value function
$V(\cdot)$ exists, that a maximum exists and it is unique;
\item Moreover, we need the \emph{envelope theorem} to derive the functional Euler equation.
\end{itemize}
\end{frame}


\begin{frame}
Example 4.5 Solving a parametric dynamic programming problem

(work from the ``guess'' that the value function has the form $V(k) = a + b\log{k}$ to obtain the decision rule: $k^{\prime} = \alpha \beta A k^{\alpha}$). 
\end{frame}



\begin{frame}
\frametitle{Simple steps to solving a dynamic optimization problem using the envelope theorem}
\begin{equation}
V(k_{t}) = \max_{c_t, k_{t+1}} \{ u(c_t) +
\beta V(k_{t+1}) \} ~, ~~ {\text{s.t.}~~} c_t = f(k_{t}) - k_{t+1}
\end{equation}
Assume the constraint binds, as before, and take FOC w.r.t $k_{t+1}$:
\begin{eqnarray} \nonumber
\frac{\partial{V}(k_t)}{\partial{k}_{t+1}} &:& \underset{{\text{Chain rule}}}{\underbrace{\frac{\partial{u}(c_t)}{\partial{c}_{t}} \frac{\partial{c_t}}{\partial{k}_{t+1}}}} + \beta\frac{\partial{V}(k_{t+1})}{\partial{k}_{t+1}} = 0 \\ \nonumber
& \therefore & u^{\prime}(c_t)(-1) + {\beta}V^{\prime}(k_{t+1}) = 0 \\ \label{FOC1}
& \therefore & u^{\prime}(c_t) = {\beta}V^{\prime}(k_{t+1})
\end{eqnarray}
We need to find $V^{\prime}(k_{t+1})$ \ldots 

\end{frame}



\begin{frame}
\frametitle{Simple steps to solving a dynamic optimization problem using the envelope theorem}

We can use the envelope theorem, by taking FOC w.r.t $k_{t}$ and then iterating forward by one period:
\begin{eqnarray} \nonumber
\frac{\partial{V}(k_t)}{\partial{k}_{t}} &:& \frac{\partial{u}(c_t)}{\partial{c}_{t}}\frac{\partial{c_t}}{\partial{f(k_t)}} \frac{\partial{f(k_t)}}{\partial{k}_t} + \beta\frac{\partial{V}(k_{t+1})}{\partial{k}_{t}} \\ \nonumber
& \therefore & u^{\prime}(c_t)(1)f^{\prime}(k_t) + \beta (0)
 \\  \label{FOC2}
 \therefore \frac{\partial{V}(k_{t+1})}{\partial{k}_{t+1}} &=& V^{\prime}(k_{t+1}) =  u^{\prime}(c_{t+1})f^{\prime}(k_{t+1}) 
\end{eqnarray}


Substitute \eqref{FOC2} into \eqref{FOC1} to get the functional Euler equation, as in our ``canonical problem'':

\begin{equation}
\underset{\text{invest ``one'' more unit}}{\underset{\text{Utility lost if you}}{\underbrace{u'(c_t)}}} 
= 
\underset{\text{next period per unit}}{\underset{\text{Utility increase}}{\underbrace{\beta{u}'(c_{t+1})}}}  \cdot  \underset{\text{invested unit}}{\underset{\text{Return on the }}{\underbrace{f'(k_{t+1})}}}
\end{equation}

\end{frame}







\section{A General Version}

\begin{frame}

\frametitle{A General Version (similiar to Krussel, Ch 4.3.)} \vskip 0.1in

In a general form, for the model economy, the social-planning
problem or, equivalently, the competitive equilibrium involves solving the following dynamic programming problem:

\begin{eqnarray}
V(x_{t},z_{t}) &=& \max_{y_t} [F(x_{t},y_{t},z_{t}) + \beta
E_{t}V(x_{t+1},z_{t+1})] \label{EqA.5}\\
s.t. ~~~~~~~ x_{t+1} &=& {\color{blue}G (x_{t},y_{t},z_{t})} \label{EqA.6}
\end{eqnarray}

where,\\

\begin{itemize}
\item $x_{t}$: a vector of state variables in t;
\item $y_{t}$: a vector of control variables in t;
\item $z_{t}$: a vector of stochastic state variables in t;
\item F($\cdot$,$\cdot$): objective function to be maximized;
\item \eqref{EqA.6}: budget constraint. We include the expectations operator because of the presence of uncertainty $z_t$.
\end{itemize}

\end{frame}

\begin{frame}

\frametitle{A General Version} \vskip 0.1in

The solution to this problem:

\begin{equation}\label{EqA.7}
y_{t} = {\color{red}H(x_{t}, z_{t})}
\end{equation}

\eqref{EqA.7} is the so called policy function (or decision rule)
which describes how the control variable behaves as a function of the
state variables in $t$.\vskip 0.1in

\pause

Since the policy function optimizes the choice of the control
variables for every permitted value of $x_{t}$, it must fulfill the
following condition:

\begin{equation}\label{EqA.8}
V (x_t, z_t) = F (x_t, {\color{red}H(x_t, z_t)}, z_t) + \beta E_t V({\color{blue}G(x_t, {\color{red}H(x_t,
z_t)}, z_t)}, z_{t+1})
\end{equation}

\end{frame}


\begin{frame}

\frametitle{A General Version} \vskip 0.1in

To find the policy function $H(x_t, z_t)$, we need the FOCs of
\eqref{EqA.5} and its envelope condition,

{\scriptsize{
\begin{eqnarray}
\frac{\partial{V}(\cdot_t)}{\partial{y_t}} &:& 0 = F_{y}(x_t, y_t, z_t) + \beta E_t [{\overset{{\color{Green}\frac{\partial{V}(x_{t+1}, z_{t+1})}{\partial{x_{t+1}}}} \frac{\partial{x}_{t+1}}{\partial{G}(\cdot)} \frac{\partial{G}(\cdot)}{\partial{y}_t} }{\overbrace{{\color{Green}{V_x({\color{blue}G(x_t, y_t,z_t)},
z_{t+1})}}G_y(x_t, y_t, z_t)}}}]  \label{EqA.9}\\ \nonumber \\ \label{EqA.9b}
\underset{{\text{\scriptsize satisfies}}}{{\color{red}H(\cdot_t)}} &:& 0 = F_{y}({\color{red}H}) + \beta E_t [{\color{Green}{V_x({\color{blue}G(H)})}}G_y({\color{red}H}))]  \\ \nonumber && {\text{But,}}~{\color{Green}V_x(\cdot_{t+1})}~{\text{is unknown}}   \\ \label{EqA.10}
{\color{Green}{\frac{\partial{V}(\cdot_{t})}{\partial{x_{t}}}}} &=& F_x  +  \beta E_t V_{x}G_x + H_x [ \underset{= 0 ~{\text{by}}~ \eqref{EqA.9b}}{\underbrace{F_{y} +  \beta E_t V_{x}G_y }}  ]
\end{eqnarray} 
}}

where,
\begin{itemize}
\item $F_{y}(x_t, y_t, z_t), F_{x}(x_t, y_t, z_t)$: vector of derivatives of the objective
function w.r.t. the control variables and state variables;
\item ${\color{Green}{V_x({\color{blue}G(x_t, y_t,z_t)},
z_{t+1})}}$: vector of derivatives of the objective
function w.r.t. the state variables in ${\color{black}{t+1}}$;
\item $G_y(x_t, y_t, z_t), G_x(x_t, y_t, z_t)$: vector of
derivatives of the budget constraints w.r.t. the control variables
and state variables.
\end{itemize}

\end{frame}

\begin{frame}
\frametitle{A General Version} \vskip 0.1in

$\Box$ Optimizing \eqref{EqA.5} such that \eqref{EqA.6} holds, implies that $G_x(\cdot)=0$, and \eqref{EqA.10}, iterated forward, becomes:
%\begin{eqnarray}
%\therefore {\color{Green}{\frac{\partial{V}(\cdot_{t+1})}{\partial{x_{t+1}}}}} &=& F_x (x_{t+1}, y_{t+1}, z_{t+1}) \\   & & + ~~  \beta E_t [V_x({\color{blue}G(x_{t+1}, y_{t+1},z_{t+1})}, z_{t+2})G_x(x_{t+1}, y_{t+1}, z_{t+1})]
%\end{eqnarray}
\begin{eqnarray}\label{EqA.11}
\therefore {\color{Green}{\frac{\partial{V}(\cdot_{t+1})}{\partial{x_{t+1}}}}} &:& V_x(x_{t+1}, z_{t+1}) = F_x (x_{t+1}, y_{t+1}, z_{t+1})
\end{eqnarray}


Therefore, the FOCs \eqref{EqA.9} give the following functional Euler equation:
\begin{equation}\label{EqA.12}
0 = F_{y}(x_t, y_t, z_t) + \beta E_t [F_x(G(x_t, y_t,z_t), y_{t+1},
z_{t+1})G_y(x_t, y_t, z_t)]
\end{equation}
Solving for $y_t$ gives the policy function \eqref{EqA.7}.
\end{frame}



\begin{frame}
\frametitle{Self-study}

Work through examples 4.1 (p. 47), 4.2 (p. 49), 4.5. (p. 63), and Ch 4.3 (pp. 67-69).

\end{frame}


%\begin{frame}
%\begin{center}
%\Huge{THE END}
%\end{center}
%\end{frame}
\end{document}
